% !TeX program = xelatex
% Вопрос 3. Только цитаты из ../tex/3/chapter.tex.
% Библиографические сведения: ../tex/3/main.tex.
% Пропуски обозначены […]. Внешние ссылки оригинала сохранены в читаемом виде.
\documentclass[12pt,a4paper]{article}
\usepackage{fontspec}
\setmainfont{Times New Roman}
\usepackage[english,russian]{babel}
\usepackage{amsmath,amssymb,cancel}
\usepackage[margin=22mm]{geometry}
\usepackage[unicode,hidelinks]{hyperref}
\DeclareMathOperator{\grad}{grad}
\DeclareMathOperator{\Div}{div}
\newcommand{\Rey}{\mathrm{Re}}
\newcommand{\topic}[1]{\par\medskip\noindent\textbf{#1}\par\nobreak\smallskip}
\setlength{\emergencystretch}{3em}
\interfootnotelinepenalty=10000
\setlength{\parindent}{0pt}
\renewcommand{\thefootnote}{\arabic{footnote})}
\newenvironment{textquote}{\begin{quote}}{\end{quote}}
\newcommand{\source}[2]{\par\nobreak\smallskip{\small Н. Е. Леонтьев. \textit{Основы теории фильтрации}. 2-е изд., 2017. Глава I, §~#1, с.~#2.}\par\medskip}
\begin{document}
% BEGIN QUESTION HEADER
\section*{Вопрос 3}
\noindent Уравнения баланса массы, импульса и энергии для течений в пористых средах. Соотношения на разрывах. Типичные граничные условия.

\medskip
\noindent\textbf{Источник:} Н. Е. Леонтьев. \textit{Основы теории фильтрации}. 2-е изд. Москва: МАКС Пресс, 2017. Глава I, §§ 3, 5--6, печатные страницы 13--17 и 24--37. Ответ составлен из дословных цитат перенесённого в TeX текста учебника.
\par\bigskip
% END QUESTION HEADER
% Цитата 1; печатные страницы: 14
\begin{textquote}
Интегральная формулировка закона сохранения массы жидкости получается из утверждения о том, что изменение массы жидкости в произвольном объеме $V$ внутри пористой среды происходит за счет притока жидкости через границу объема $\Sigma$:
\[
\frac{d}{dt}\int_V\rho m\,dV=-\int_\Sigma\rho\boldsymbol{u}\cdot\boldsymbol{n}\,d\sigma.
\]
Здесь $\boldsymbol{n}$ --- вектор внешней нормали к поверхности $\Sigma$ (из-за такого выбора нормали интеграл в правой части стоит со знаком минус), а объем $V$ неподвижен относительно пористой среды. Если входящие в это уравнение функции являются достаточно гладкими, то в левой части можно занести производную под знак интеграла, а правую часть преобразовать с помощью теоремы Гаусса--Остроградского:
\[
\int_V\frac{\partial(\rho m)}{\partial t}\,dV=-\int_V\Div(\rho\boldsymbol{u})\,dV,
\]
откуда в силу произвольности объема $V$ получается дифференциальная запись закона сохранения массы --- \textit{уравнение неразрывности}:
\begin{equation}
\frac{\partial(\rho m)}{\partial t}+\Div(\rho\boldsymbol{u})=0.
\tag{3.1}\label{q3:eq-3-1}
\end{equation}
Если плотность жидкости постоянна, а пористая среда не деформируется (т.е. пористость зависит только от координат), то первое слагаемое обращается в ноль и уравнение неразрывности приобретает вид
\begin{equation}
\Div\boldsymbol{u}=0.
\tag{3.2}\label{q3:eq-3-2}
\end{equation}
\source{3}{14}
\end{textquote}

% Цитата 2; печатные страницы: 14--15
\begin{textquote}
Исторически уравнение баланса импульса было экспериментально получено французским гидравликом Анр\'и Дарс\'и (H.Darcy, 1856) для случая медленного стационарного движения несжимаемой жидкости в неподвижной изотропной пористой
среде. В современных обозначениях это соотношение, называемое \textit{законом Дарси}, для фильтрации в поле тяжести имеет вид
\begin{equation}
0=-\grad p+\rho\boldsymbol{g}-\frac{\mu}{k}\boldsymbol{u},
\tag{3.3}\label{q3:eq-3-3}
\end{equation}
где $\boldsymbol{g}$ --- ускорение свободного падения, $\mu$ --- коэффициент динамической вязкости жидкости, $k$ --- коэффициент, называемый проницаемостью пористой среды, который зависит от типа пористой среды.
\source{3}{14--15}
\end{textquote}

% Цитата 3; печатные страницы: 25
\begin{textquote}
В случае \textit{недеформируемого неподвижного} пористого скелета, которым мы для простоты ограничимся, при вычислении среднего от производных справедливы равенства
\begin{equation}
\left\langle\frac{\partial A}{\partial t}\right\rangle=\frac{\partial\langle A\rangle}{\partial t},\qquad
\left\langle\frac{\partial A}{\partial x^i}\right\rangle=\frac{\partial\langle A\rangle}{\partial x^i}+\frac{1}{|V|}\int_{\Sigma_{\text{ж-с}}}An_i\,d\sigma,
\tag{5.1}\label{q3:eq-5-1}
\end{equation}
где $\Sigma_{\text{ж-с}}$ --- граница жидкости и скелета, попадающая внутрь $V$, $\boldsymbol{n}=(n_x;n_y;n_z)$ --- вектор внешней нормали к заштрихованному на рис.~7 объему $V_{\text{ж}}$ жидкости, находящейся внутри сферы.
\source{5}{25}
\end{textquote}

% Цитата 4; печатные страницы: 26--27
\begin{textquote}
Теперь с помощью (5.1) несложно получить осредненные уравнения фильтрации. Проделаем это для случая однородной \textit{несжимаемой} ньютоновской жидкости.

\topic{Осреднение микроскопического уравнения неразрывности}
Осреднение уравнения неразрывности $\Div\boldsymbol{v}=0$ производится очень просто:
\[
\langle\nabla_iv^i\rangle=\nabla_i\langle v^i\rangle+\frac{1}{|V|}\int_{\Sigma_{\text{ж-с}}}v^in_i\,d\sigma=0,
\]
откуда с учетом условия прилипания получаем выражение закона сохранения массы в терминах средних параметров:
\begin{equation}
\nabla_i\langle v^i\rangle\equiv\Div\langle\boldsymbol{v}\rangle=0.
\tag{5.2}\label{q3:eq-5-2}
\end{equation}
Величина, стоящая под знаком дивергенции, есть не что иное, как скорость фильтрации:
\[
\langle\boldsymbol{v}\rangle=\boldsymbol{u},
\]
так что (5.2) совпадает с (3.2).

\topic{Осреднение микроскопического уравнения движения}
Для осреднения удобно взять уравнение движения в дивергентном виде
\[
\rho\left(\frac{\partial v^i}{\partial t}+\nabla_k(v^kv^i)\right)=\nabla_kp^{ik}+\rho F^i,
\]
в котором $p^{ik}$ --- компоненты тензора напряжений, $F^i$ --- компоненты вектора массовой плотности внешних объемных сил. Применяя (5.1) и учитывая условие прилипания, получим
\[
\rho\left(\frac{\partial\langle v^i\rangle}{\partial t}+\nabla_k\langle v^kv^i\rangle\right)=\nabla_k\langle p^{ik}\rangle+F^i_{\text{с-ж}}+\rho\langle F^i\rangle,
\]
\[
F^i_{\text{с-ж}}=\frac{1}{|V|}\int_{\Sigma_{\text{ж-с}}}p^{ik}n_k\,d\sigma,
\]
где сила $\boldsymbol{F}_{\text{с-ж}}$, вид которой будет обсуждаться чуть позже, представляет собой \textit{полную} силу (в расчете на единицу объема), действующую со стороны скелета на жидкость внутри объема $V$, т.е. включает в себя как результат действия вязких касательных напряжений, так и результат суммирования нормальных напряжений на границе жидкости и скелета. Первое слагаемое в правой части, снова с учетом условия прилипания, принимает вид
\[
\begin{aligned}
\nabla_k\langle p^{ik}\rangle
&=\nabla_k\left\langle-pg^{ik}+\mu(\nabla^iv^k+\nabla^kv^i)\right\rangle=\\
&=-\nabla^i\langle p\rangle+\cancel{\mu\nabla_k\nabla^i\langle v^k\rangle}+\mu\nabla_k\nabla^k\langle v^i\rangle,
\end{aligned}
\]
где зачеркнутый член пропадает вследствие (5.2), что в результате дает
\begin{equation}
\begin{aligned}
\rho\left(\frac{\partial\langle v^i\rangle}{\partial t}+\nabla_k\langle v^kv^i\rangle\right)
&=-\nabla^i\langle p\rangle+\mu\triangle\langle v^i\rangle\\
&\quad+\rho\langle F^i\rangle+F^i_{\text{с-ж}}.
\end{aligned}
\tag{5.3}\label{q3:eq-5-3}
\end{equation}
Уравнение (5.3) является точным, и дальнейшее его упрощение связанно с дополнительными гипотезами и оценкой входящих в него членов.
\source{5}{26--27}
\end{textquote}

% Цитата 5; печатные страницы: 28
\begin{textquote}
Удобно выделить в явном виде \textit{силу сопротивления} (в расчете \textit{на единицу объема жидкости}), связанную --- прямо или косвенно --- с действием вязких сил и обращающуюся в ноль при $\boldsymbol{u}=0$, и переписать уравнение движения в виде
\begin{equation}
\begin{gathered}
0=-\grad p+\rho\boldsymbol{F}+\boldsymbol{F}_{\text{сопр}},\\
\boldsymbol{F}_{\text{сопр}}=\frac{1}{m}(\boldsymbol{F}_{\text{с-ж}}-p\grad m).
\end{gathered}
\tag{5.4}\label{q3:eq-5-4}
\end{equation}
\source{5}{28}
\end{textquote}

% Цитата 6; печатные страницы: 29
\begin{textquote}
В случае постоянной пористости закон (5.4) должен совпадать с законом Дарси, поэтому сила сопротивления
\[
\boldsymbol{F}_{\text{сопр}}=\frac{1}{m}\boldsymbol{F}_{\text{с-ж}}=-\frac{\mu}{k}\boldsymbol{u}.
\]
\source{5}{29}
\end{textquote}

% Цитата 7; печатные страницы: 29--30
\begin{textquote}
Осреднение уравнения энергии для жидкости
\[
\begin{aligned}
\frac{\partial}{\partial t}\left(\rho\left(U+\frac{v^2}{2}\right)\right)
&+\nabla_k\left(\rho v^k\left(U+\frac{v^2}{2}\right)\right)=\\
&=\rho F_kv^k+\nabla_k(p^{ik}v_i)-\nabla_kq^k+\rho Q
\end{aligned}
\]
и уравнения энергии для скелета
\[
\frac{\partial}{\partial t}(\rho_{\text{с}}U_{\text{с}})=-\nabla_kq_{\text{с}}^k+\rho_{\text{с}}Q_{\text{с}},
\]
где $U$ и $Q$ --- массовые плотности соответственно внутренней энергии и объемного притока энергии, $\boldsymbol{q}$ --- вектор потока тепла, а индекс «с» обозначает величины, относящиеся к скелету, дает, с учетом условия прилипания, уравнения
\[
\begin{aligned}
\frac{\partial}{\partial t}\left\langle\rho\left(U+\frac{v^2}{2}\right)\right\rangle
&+\nabla_k\left\langle\rho v^k\left(U+\frac{v^2}{2}\right)\right\rangle=\\
&=\langle\rho F_kv^k\rangle+\nabla_k\langle p^{ik}v_i\rangle-\nabla_k\langle q^k\rangle\\
&\quad-\frac{1}{|V|}\int_{\Sigma_{\text{ж-с}}}q^kn_k\,d\sigma+\rho\langle Q\rangle,
\end{aligned}
\]
\[
\frac{\partial}{\partial t}\langle\rho_{\text{с}}U_{\text{с}}\rangle_{\text{с}}
=-\nabla_k\langle q_{\text{с}}^k\rangle_{\text{с}}+\frac{1}{|V|}\int_{\Sigma_{\text{ж-с}}}q_{\text{с}}^kn_k\,d\sigma+\rho_{\text{с}}\langle Q_{\text{с}}\rangle_{\text{с}},
\]
при сложении которых интегральные члены сокращаются (в силу условия $(\boldsymbol{q}-\boldsymbol{q}_{\text{с}})\cdot\boldsymbol{n}=0$ сохранения потока энергии на контактном разрыве --- границе жидкости и скелета), в результате чего получается требуемое уравнение энергии
\begin{equation}
\begin{aligned}
\frac{\partial}{\partial t}\left\{\left\langle\rho\left(U+\frac{v^2}{2}\right)\right\rangle+\langle\rho_{\text{с}}U_{\text{с}}\rangle_{\text{с}}\right\}
&+\nabla_k\left\langle\rho v^k\left(U+\frac{v^2}{2}\right)\right\rangle=\\
&=\langle\rho F_kv^k\rangle+\nabla_k\langle p^{ik}v_i\rangle\\
&\quad-\nabla_k\bigl(\langle q^k\rangle+\langle q_{\text{с}}^k\rangle_{\text{с}}\bigr)\\
&\quad+\rho\langle Q\rangle+\rho_{\text{с}}\langle Q_{\text{с}}\rangle_{\text{с}}.
\end{aligned}
\tag{5.6}\label{q3:eq-5-6}
\end{equation}
\source{5}{29--30}
\end{textquote}

% Цитата 8; печатные страницы: 30--31
\begin{textquote}
Этому уравнению можно придать более простой вид, если сделать следующие дополнительные предположения, выполняющиеся во многих практических приложениях.
\begin{enumerate}\sloppy
\item[а)] Кинетической энергией жидкости при движении внутри пор можно пренебречь по сравнению со внутренней энергией (в промысловых условиях $U\sim10^4$--$10^6$~м$^2$/с$^2$ $\gg v^2/2$).
\item[б)] Объемные притоки энергии отсутствуют: $Q=Q_{\text{с}}=0$.
\item[в)] Температуры жидкости и скелета можно считать одинаковыми, так как характерное время выравнивания их температур $t_*\sim\rho cd^2/\varkappa$, где $c$ и $\varkappa$ --- коэффициенты теплоемкости и теплопроводности, мало по сравнению с характерным временем фильтрационных процессов (величина $t_*$ для песка с $d\sim10^{-2}$~см порядка $10^{-1}$~с).
\item[г)] Давление и температура являются медленно меняющимися функциями пространственных координат, в то время как скорость жидкости внутри пор является быстро меняющейся функцией. Это позволяет приближенно выполнять осреднение произведений:
\[
\langle pv^k\rangle\approx p\langle v^k\rangle=pu^k,\quad
\langle\rho Uv^k\rangle\approx\rho U\langle v^k\rangle=\rho Uu^k
\]
(в правых частях равенств стоят средние значения $p$, $U$ \textit{по жидкой фазе}).
\item[д)] В слагаемом, содержащем компоненты тензора напряжений, при осреднении можно пренебречь тензором вязких напряжений $\tau^{ij}$:
\[
\begin{aligned}
\nabla_k\langle p^{ik}v_i\rangle
&=-\nabla_k\langle pv^k\rangle+\nabla_k\langle\tau^{ik}v_i\rangle\\
&\approx-\nabla_k\bigl(p\langle v^k\rangle\bigr)=-\Div(p\boldsymbol{u}).
\end{aligned}
\]
В самом деле, $\tau^{ij}\sim\mu u/d$, а порядок изменения давления $\delta p$ (в масштабах задачи $L$) оценивается с помощью закона Дарси: $\delta p/L\sim\mu u/d^2$, откуда $\tau^{ij}/\delta p\sim d/L\ll1$. С физической точки зрения это означает, что мощность сил давления много больше мощности вязких сил при деформации жидкости.
\item[е)] Для суммарного вектора потока тепла справедлив закон теплопроводности Фурье с некоторым эффективным коэффициентом теплопроводности $\varkappa_{\text{эф}}$:
\begin{equation}
\boldsymbol{q}_{\text{сум}}\equiv\langle\boldsymbol{q}\rangle+\langle\boldsymbol{q}_{\text{с}}\rangle_{\text{с}}=-\varkappa_{\text{эф}}\grad T.
\tag{5.7}\label{q3:eq-5-7}
\end{equation}
Обратим внимание, что $\varkappa_{\text{эф}}$ в общем случае не выражается только через пористость $m$ и коэффициенты теплопроводности жидкости $\varkappa$ и скелета $\varkappa_{\text{с}}$, а зависит от внутренней структуры пористой среды [4].
\end{enumerate}
После упрощений, проведенных с применением перечисленных предположений, уравнение энергии (5.6) приобретает вид
\begin{equation}
\begin{aligned}
\frac{\partial}{\partial t}(\rho Um+\rho_{\text{с}}U_{\text{с}}(1-m))+\Div(\rho U\boldsymbol{u})&=\\
&=\rho\boldsymbol{F}\cdot\boldsymbol{u}-\Div(p\boldsymbol{u})-\Div\boldsymbol{q}_{\text{сум}}.
\end{aligned}
\tag{5.8}\label{q3:eq-5-8}
\end{equation}
Отметим, что учет сжимаемости среды принципиально не изменит проводившихся при выводе (5.8) рассуждений, поэтому это уравнение может использоваться и для описания тепловых явлений при фильтрации газа.
\source{5}{30--31}
\end{textquote}

% Цитата 9; печатные страницы: 32
\begin{textquote}
Связь параметров при переходе через поверхность разрыва получается, как всегда, из интегральных законов сохранения, которые при крупномасштабном описании имеют вид
\begin{equation}
\begin{gathered}
\frac{d}{dt}\int_V\rho m\,dV=-\int_\Sigma\rho u_n\,d\sigma,\\
0=-\int_\Sigma mp\boldsymbol{n}\,d\sigma+\int_V m\left(\rho\boldsymbol{F}-\frac{\mu}{k}\boldsymbol{u}\right)\,dV,\\
\begin{aligned}
\frac{d}{dt}\int_V(\rho Um+\rho_{\text{с}}U_{\text{с}}(1-m))\,dV&=\\
&=-\int_\Sigma\bigl((\rho U+p)u_n+\boldsymbol{q}_{\text{сум}}\cdot\boldsymbol{n}\bigr)\,d\sigma\\
&\quad+\int_V(\rho\boldsymbol{F}\cdot\boldsymbol{u}+\rho Qm+\rho_{\text{с}}Q_{\text{с}}(1-m))\,dV
\end{aligned}
\end{gathered}
\tag{6.1}\label{q3:eq-6-1}
\end{equation}
для любого \textit{неподвижного относительно пористой среды}\footnote[1]{Пористая среда предполагается неподвижной относительно некоторой инерциальной системы отсчета.} объема $V$ с границей $\Sigma$. Все интегральные соотношения имеют общую структуру
\[
\frac{d}{dt}\int_V A\,dV=-\int_\Sigma\boldsymbol{B}\cdot\boldsymbol{n}\,d\sigma+\int_V C\,dV,
\]
которой соответствует условие на разрыве вида
\begin{equation}
[AD-\boldsymbol{B}\cdot\boldsymbol{n}]=I,
\tag{6.2}\label{q3:eq-6-2}
\end{equation}
где, как обычно, квадратными скобками обозначается скачок величины в скобках при переходе через разрыв, $D$ --- скорость движения поверхности разрыва по нормали к ней, $I$ --- поверхностный приток, соответствующий величине $A$.
\source{6}{32}
\end{textquote}

% Цитата 10; печатные страницы: 33
\begin{textquote}
Применение (6.2) к законам (6.1) дает условия сохранения массы, импульса и энергии
\begin{equation}
\begin{gathered}
\relax[\rho(mD-u_n)]=M,\qquad[-mp\boldsymbol{n}]=\boldsymbol{R},\\
[(\rho Um+\rho_{\text{с}}U_{\text{с}}(1-m))D-(\rho U+p)u_n-\boldsymbol{q}_{\text{сум}}\cdot\boldsymbol{n}]=E,
\end{gathered}
\tag{6.3}\label{q3:eq-6-3}
\end{equation}

где $M$, $\boldsymbol{R}$ и $E$ --- дополнительные поверхностные притоки массы, импульса и энергии на разрыве, которые, вообще говоря, отличны от нуля. Например, при фильтрации воды через границу двух пористых сред с разными пористостями возникает дополнительная поверхностная сила, действующая со стороны скелета на жидкость\footnote[1]{Природу этой силы легко понять на примере модельного течения через стык двух капилляров с разными диаметрами.}, равная по модулю $|p[m]|$ (при изотермической фильтрации из условия сохранения энергии в этой задаче следует непрерывность давления). Другим примером поверхностной силы может служить возникающая на границе двух жидкостей, находящихся в пористой среде, капиллярная разность давлений, величина которой имеет порядок $R\sim m\sigma/d$, где $\sigma$ --- коэффициент поверхностного натяжения.
\source{6}{33}
\end{textquote}

% Цитата 11; печатные страницы: 33
\begin{textquote}
На границе пористой среды с непроницаемой стенкой или с другой пористой средой, проницаемостью которой можно пренебречь, ставится \textit{условие непротекания} $\boldsymbol{u}\cdot\boldsymbol{n}=0$.
\source{6}{33}
\end{textquote}

% Цитата 12; печатные страницы: 33--34
\begin{textquote}
Обратим внимание, что при микроскопическом описании движения жидкости на поверхности пор ставится условие прилипания, т.е. требуется обращение в ноль как нормальной, так и касательной составляющей скорости. При осредненном описании на границе можно поставить условие только на объемный поток
жидкости через границу --- он должен быть равен нулю, т.к. стенка непроницаема. Касательная же составляющая скорости фильтрации $\boldsymbol{u}_\tau$ может быть отлична от нуля --- жидкость может мигрировать в пристеночном слое вдоль границы. С математической точки зрения эта разница в граничных условиях связана с тем, что микроскопическое уравнение движения (уравнение Навье--Стокса) содержит вторые производные от скорости по координатам, а макроскопическое уравнение движения (закон Дарси) зависит от скорости алгебраически.
\source{6}{33--34}
\end{textquote}

% Цитата 13; печатные страницы: 34
\begin{textquote}
Во многих задачах встречаются границы пористой среды, насыщенной жидкостью, со внешним объемом той же самой жидкости либо границы пористой среды, в которой фильтруется газ, с тем же газом в свободном состоянии.

[…]

Как правило, на таких границах ставятся условия сохранения массы $[\rho u_n]=0$ и непрерывности давления $[p]=0$ (в задачах изотермического движения несжимаемой жидкости последнее условие следует из закона сохранения энергии).

В водоемах распределение давления с глубиной практически не зависит от наличия течений, поэтому в гидротехнических задачах внешнее давление, как правило, определяется по гидростатическому закону $p=p_{\text{атм}}+\rho gh$, где $h$ --- вертикальное расстояние от рассматриваемой точки на дне до поверхности жидкости. В задачах фильтрации газа обычно можно считать внешние давление и плотность газа заданными постоянными величинами или функциями времени.
\source{6}{34}
\end{textquote}

% Цитата 14; печатные страницы: 35
\begin{textquote}
Если разные части пористой среды насыщены двумя разными жидкостями или жидкостью и газом, то граница между этими областями является поверхностью разрыва\footnote[1]{Подчеркнем, что здесь не рассматривается многофазная фильтрация (течения смесей, эмульсий, жидкостей с пузырьками газа и т.п. в пористой среде), когда характерные размеры связных областей одной фазы меньше или порядка характерного размера пор.}.

[…]

При изотермической фильтрации на подвижной границе ставятся обычно общее условие сохранения массы $[\rho(mD-u_n)]=0$ и (в пренебрежении капиллярными силами) условие непрерывности давления $[p]=0$.
\source{6}{35}
\end{textquote}

% Цитата 15; печатные страницы: 36
\begin{textquote}
В ряде гидротехнических задач на части границы пористой среды с атмосферой жидкость может вытекать из пористой среды и, скажем, стекать вдоль границы или испаряться. Такие участки называются \textit{поверхностями высачивания} и встречаются, например, на стенках земляных плотин, через которые фильтруется вода, или на стенках колодцев, вырытых в водонасыщенном грунте (рис.~9).

На поверхности высачивания ставятся условия непрерывности давления (в пренебрежении капиллярными эффектами давление жидкости в пористой среде на границе совпадает с атмосферным) и положительности проекции скорости фильтрации на внешнюю нормаль $\boldsymbol{u}\cdot\boldsymbol{n}>0$ (жидкость может только вытекать из пористой среды). Сама же величина скорости фильтрации на границе до решения задачи неизвестна.
\source{6}{36}
\end{textquote}

% Цитата 16; печатные страницы: 36--37
\begin{textquote}
\topic{Граница двух пористых сред}
Примерами подобных границ являются поверхности раздела между расположенными рядом природными пластами с разными проницаемостями или, скажем, граница между фильтром и внешней пористой осадочной коркой, откладывающейся на его внешней поверхности при его загрязнении.


При фильтрации несжимаемой жидкости на такой границе ставятся условия непрерывности нормальной компоненты скорости фильтрации $[u_n]=0$ и, как уже говорилось, непрерывности давления $[p]=0$. Отметим, что \textit{касательные} составляющие скорости фильтрации $\boldsymbol{u}_\tau$ по разные стороны от разрыва в общем случае отличаются (рис.~10). В самом деле, в силу условия непрерывности давления величина напора $H$ непрерывна на разрыве, поэтому проектирование закона Дарси (3.6), записанного по обе стороны поверхности, на касательную плоскость к поверхности дает условие
\[
\frac{|\boldsymbol{u}_{1\tau}|}{C_1}=\frac{|\boldsymbol{u}_{2\tau}|}{C_2},
\]
т.е. касательные составляющие скорости фильтрации относятся как коэффициенты фильтрации. Отсюда, с учетом условия
$[u_n]=0$, получается соотношение на углы $\alpha_1$ и $\alpha_2$, которые образуют векторы скорости фильтрации с вектором нормали:
\[
\frac{\operatorname{tg}\alpha_1}{C_1}=\frac{\operatorname{tg}\alpha_2}{C_2}.
\]
Подчеркнем, что в общем случае (например, при фильтрации газа через границу двух сред с разными пористостями, когда внутрипоровые скорости движения частиц газа велики) скачок давления на границе может быть отличен от нуля и должен определяться из дополнительных соотношений\footnote[1]{Аналогичная ситуация встречается, например, при течении газа в трубе с внезапным расширением или при протекании газа через перфорированные перегородки.}.
\source{6}{36--37}
\end{textquote}

\end{document}
